\chapter{并发再探}
\label{CH:LOCK2}

在内核设计中，同时获得良好的并行性能、保证并发下的正确性以及编写可理解的代码是一个巨大挑战。
直接使用锁是实现正确性的最佳途径，但并不总是可行的。
本章重点介绍 xv6 中被迫以复杂方式使用锁的例子，以及 xv6 使用类似锁的技术而非实际锁的例子。

\section{锁的模式}

缓存项的锁管理通常是个难题。
例如，文件系统的块缓存 \lineref{kernel/bio.c:/^struct/} 存储了最多 {\tt NBUF} 个磁盘块的副本。
确保每个磁盘块在缓存中至多有一个副本至关重要；否则，不同进程可能对本应是同一块的不同副本做出相互冲突的修改。
每个缓存的块存储在一个 {\tt struct buf} \fileref{kernel/buf.h} 结构体中。
{\tt struct buf} 有一个锁字段，用于帮助确保同一时间只有一个进程使用给定的磁盘块。
然而，这个锁还不够：如果一个块根本不在缓存中，而两个进程同时想要使用它，会发生什么？
此时没有对应的 {\tt struct buf}（因为该块尚未被缓存），因此也就没有任何东西可以被锁。
xv6 通过为缓存块的标识集合关联一个额外的锁（{\tt bcache.lock}）来处理这种情况。
需要检查某个块是否已被缓存（例如 {\tt bget} \lineref{kernel/bio.c:/^bget/}）或修改缓存块集合的代码必须持有 {\tt bcache.lock}；
在该代码找到所需的块和 {\tt struct buf} 之后，它可以释放 {\tt bcache.lock}，然后仅锁定特定的块。
这是一种常见的模式：一个锁用于保护项集合，再加上每个项自身的一个锁。

通常情况下，获取锁的函数会释放该锁。但更精确的看法是：锁在需要表现为原子性的操作序列开始时被获取，在该序列结束时被释放。
如果该序列在不同的函数、不同的线程或不同的 CPU 上开始和结束，那么锁的获取和释放也必须跨这些执行环境。
锁的功能是强制其他使用者等待，而不是将数据固定给某个特定的执行代理。
一个例子是 {\tt yield} \lineref{kernel/proc.c:/^yield/} 中的 {\tt acquire}，它由调度器线程释放，而不是由获取锁的进程释放。
另一个例子是 {\tt ilock} \lineref{kernel/fs.c:/^ilock/} 中的 {\tt acquiresleep}；这段代码在读取磁盘时常会睡眠；它可能在不同的 CPU 上唤醒，这意味着锁可能在不同的 CPU 上被获取和释放。

释放一个受其内嵌锁保护的对象是一件微妙的事情，因为仅仅拥有该锁并不足以保证释放操作是正确的。
问题出现在当其他线程正在 {\tt acquire} 中等待使用该对象时；释放该对象会隐式地释放其内嵌的锁，这会导致等待的线程发生故障。
一种解决方案是跟踪该对象的引用计数，使其仅在最后一个引用消失时才被释放。
参见 {\tt pipeclose} \lineref{kernel/pipe.c:/^pipeclose/} 的例子；
{\tt pi->readopen} 和 {\tt pi->writeopen} 跟踪指向该管道的文件描述符是否还存在。

通常我们会看到锁用于保护对一组相关项进行读写的操作序列；锁确保其他线程只能看到完整的更新序列（只要它们也进行加锁）。
那么，当更新只是对单个共享变量的简单写入时，情况如何呢？
例如，\texttt{setkilled} 和 \texttt{killed} \lineref{kernel/proc.c:/^setkilled/} 在对 \lstinline{p->killed} 的简单使用中使用了锁。
如果没有锁，一个线程可能在另一个线程读取 \lstinline{p->killed} 的同时写入该变量。
这是一种\indextextx{竞争}，而 C 语言标准规定竞争会导致\indextext{未定义行为}，这意味着程序可能崩溃或产生错误的结果\footnote{“线程与数据竞争” 见 \url{https://en.cppreference.com/w/c/language/memory_model}}。
锁能防止竞争并避免未定义行为。

竞争会破坏程序的一个原因是，如果没有锁或等效的构造，编译器生成的机器代码可能以与原始 C 代码截然不同的方式读写内存。
例如，调用 \texttt{killed} 的线程的机器代码可能会将 \lstinline{p->killed} 复制到一个寄存器中，并且只读取这个缓存的值；这意味着该线程可能永远看不到对 \lstinline{p->killed} 的任何写入。
锁能防止此类缓存行为。

% sleep locks.
% hand-over-hand locking in namei.
%   namei's use of refcount to prevent changing underfoot,
%     and lock when actually using it.
% example of where deadlock is avoided? namei?
% spawn a thread to evade a lock order problem?

\section{类似锁的模式}

在许多地方，xv6 以类似锁的方式使用引用计数或标志，来表示某个对象已被分配且不应被释放或重用于其他用途。
进程的 {\tt p->state} 就以这种方式工作，{\tt file}、{\tt inode} 和 {\tt buf} 结构体中的引用计数也是如此。
虽然在这些情况下，锁用于保护标志或引用计数，但真正防止对象被过早释放的是引用计数机制本身。

文件系统使用 {\tt struct inode} 引用计数作为一种可以被多个进程持有的共享锁，以避免使用普通锁可能导致的死锁。
例如，{\tt namex} \lineref{kernel/fs.c:/^namex/} 中的循环依次锁定每个路径名分量所指代的目录。
然而，{\tt namex} 必须在循环结束时释放每个锁，因为如果它持有多个锁，当路径名中包含点（例如 {\tt a/./b}）时，可能会与自身发生死锁。
它也可能与涉及该目录和 {\tt ..} 的并发查找操作发生死锁。
如第~\ref{CH:FS} 章所述，解决方案是让循环将目录 inode 的引用计数增加后带入下一次迭代，但不对其加锁。

某些数据项在不同时间受不同机制的保护，有时可能并非通过显式锁，而是隐式地受 xv6 代码结构的保护以避免并发访问。
例如，当一个物理页空闲时，它受到 \texttt{kmem.lock} \lineref{kernel/kalloc.c:/^. kmem;/} 的保护。
如果该页随后被分配用作管道 \lineref{kernel/pipe.c:/^pipealloc/}，它则会受到一个不同的锁（内嵌的 \lstinline{pi->lock}）的保护。
如果该页被重新分配给新进程的用户内存，则根本不受任何锁的保护。
相反，分配器不会将该页再分配给任何其他进程（直到它被释放）这一事实，保护了它免受并发访问。
新进程内存的所有权是复杂的：首先父进程在 {\tt fork} 中分配并操作它，然后子进程使用它，之后（子进程退出后）父进程再次拥有该内存并将其传递给 {\tt kfree}。
这里有两个教训：一个数据对象在其生命周期的不同阶段可能以不同的方式受到并发保护，并且这种保护可能采取隐式结构而非显式锁的形式。

最后一个类似锁的例子是，在调用 {\tt mycpu()} \lineref{kernel/proc.c:/^myproc/} 时需要禁用中断。
禁用中断会使调用代码相对于可能导致上下文切换的定时器中断具有原子性，从而防止进程被移动到不同的 CPU。

\section{完全无锁}

有几个地方 xv6 在没有锁的情况下共享可变数据。一个是在自旋锁的实现中，尽管可以将 RISC-V 的原子指令视为依赖于硬件实现的锁。另一个是 {\tt main.c} \lineref{kernel/main.c:/^volatile/} 中的 {\tt started} 变量，用于防止其他 CPU 在零号 CPU 完成 xv6 初始化之前运行；{\tt volatile} 确保编译器实际生成加载和存储指令。

xv6 中存在这样的情况：一个 CPU 或线程写入某些数据，而另一个 CPU 或线程读取该数据，但没有特定的锁专门用于保护这些数据。例如，在 {\tt fork} 中，父进程写入子进程的用户内存页，而子进程（可能在不同 CPU 上的不同线程）读取这些页；没有锁显式保护这些页。严格来说，这不是一个锁问题，因为子进程在父进程完成写入之后才开始执行。但这可能存在内存顺序问题（参见第~\ref{CH:LOCK} 章），因为如果没有内存屏障，没有理由认为一个 CPU 能看到另一个 CPU 的写入。然而，由于父进程释放锁，而子进程在启动时获取锁，{\tt acquire} 和 {\tt release} 中的内存屏障确保了子进程的 CPU 能看到父进程的写入。

\section{并行性}

加锁主要是为了抑制并行性以保证正确性。由于性能也很重要，内核设计者常常需要思考如何以既能保证正确性又能允许并行性的方式使用锁。虽然 xv6 并非为高性能而系统化设计，但仍然值得思考 xv6 的哪些操作可以并行执行，以及哪些操作可能会在锁上发生冲突。

xv6 中的管道是一个具有相当好并行性的例子。每个管道都有自己的锁，因此不同进程可以在不同 CPU 上并行地读写不同的管道。然而，对于给定的管道，写入者和读取者必须等待对方释放锁；它们不能同时读写同一个管道。另外，从空管道读取（或向满管道写入）必须阻塞，但这并非由锁机制导致。

上下文切换是一个更复杂的例子。两个分别在其各自 CPU 上执行的内核线程可以同时调用 {\tt yield}、{\tt sched} 和 {\tt swtch}，并且这些调用将并行执行。每个线程都持有一个锁，但它们是不同的锁，因此不必相互等待。然而，一旦进入 {\tt scheduler}，两个 CPU 在搜索进程表以查找处于 {\tt RUNNABLE} 状态的进程时，可能会在锁上发生冲突。也就是说，xv6 在上下文切换期间很可能从多个 CPU 获得性能提升，但可能没有达到最大可能的程度。

另一个例子是不同 CPU 上的不同进程并发调用 {\tt fork}。这些调用可能因为 {\tt pid\_lock}、{\tt kmem.lock} 以及搜索进程表以查找 {\tt UNUSED} 进程所需的每进程锁而相互等待。另一方面，两个正在 fork 的进程可以完全并行地复制用户内存页和格式化页表页。

上述每个例子中的锁机制都在某些情况下牺牲了并行性能。在每种情况下，都有可能通过更精细的设计来获得更多的并行性。是否值得这样做取决于细节：相关操作被调用的频率、代码在争用锁上花费的时间、可能同时运行冲突操作的有多少 CPU、代码的其他部分是否为更严格的瓶颈。通常很难猜测给定的锁方案是否会导致性能问题，或者新设计是否显著更好，因此通常需要在真实负载上进行测量。

\section{练习}

\begin{enumerate}

\item 修改 xv6 的管道实现，允许对同一管道的读和写在不同 CPU 上并行进行。

\item 修改 xv6 的 \texttt{scheduler()}，以减少不同 CPU 同时寻找可运行进程时的锁争用。

\item 消除 xv6 的 \texttt{fork()} 中的一些串行化操作。

\end{enumerate}
